\documentclass{article}
\usepackage[T1]{fontenc}
\usepackage{lmodern}
\usepackage{graphicx}
\usepackage{amsmath,amssymb}
\usepackage[a4paper,margin=2cm]{geometry}
\usepackage[absolute,overlay]{textpos}
\setlength{\TPHorizModule}{1mm}
\setlength{\TPVertModule}{1mm}
\usepackage[indonesian]{babel}
\usepackage{hyperref}
\setlength{\emergencystretch}{2em}
% unit: Halaman 18

\begin{document}

% === Halaman 18 (llm) ===

\begin{itemize}
    \item Penemu algoritma $RSA$ pada tahun 1976 menyarankan nilai $p$ dan $q$ panjangnya lebih dari 100 digit. Dengan demikian hasil kali $n = p \times q$ akan berukuran lebih dari 200 digit.
    \item Usaha untuk mencari faktor bilangan 200 digit membutuhkan waktu komputasi selama 4 milyar tahun, sedangkan untuk bilangan 500 digit membutuhkan waktu $10^{25}$ tahun (dengan asumsi bahwa algoritma pemfaktoran yang digunakan adalah algoritma yang tercepat saat ini dan komputer yang dipakai mempunyai kecepatan 1 milidetik).
    \item Algoritma pemfaktoran yang tercepat saat ini memiliki kompleksitas
    \[ O(\exp(\sqrt[3]{\frac{64}{9}}b(\log(b)^2))) \]
    untuk bilangan bulat $n$ sepanjang b-bit.
\end{itemize}

\end{document}
